\documentclass[12pt,a4paper,oneside]{report}

% --- Idioma, tipografía y página ---
\usepackage[utf8]{inputenc}
\usepackage[T1]{fontenc}
\usepackage[spanish,es-nodecimaldot]{babel}
\usepackage{lmodern}
\usepackage{microtype}
\usepackage[a4paper,margin=2.7cm,headheight=16pt]{geometry}

% --- Tablas, figuras y diagramas ---
\usepackage{graphicx}
\usepackage{float}
\usepackage{booktabs,longtable,array,tabularx}
\usepackage{enumitem}
\usepackage{xcolor}
\usepackage{listings}
\usepackage{tikz}
\usetikzlibrary{arrows.meta,positioning,shapes.geometric,fit}
\usepackage{caption}

% --- Estilo editorial ---
\definecolor{azulNoche}{HTML}{17324D}
\definecolor{azulMedio}{HTML}{286A8A}
\definecolor{turquesa}{HTML}{2A9D8F}
\definecolor{grisTexto}{HTML}{4B5563}
\definecolor{grisFondo}{HTML}{F3F6F8}
\definecolor{dorado}{HTML}{D5A84B}
\usepackage{titlesec}
\titleformat{\chapter}[display]
  {\normalfont\bfseries\color{azulNoche}}
  {\filleft\Large\chaptertitlename\ \thechapter}
  {0.6ex}{\titlerule\vspace{1ex}\Huge}
\titleformat{\section}{\Large\bfseries\color{azulMedio}}{\thesection}{0.7em}{}
\titleformat{\subsection}{\large\bfseries\color{azulNoche}}{\thesubsection}{0.7em}{}
\titlespacing*{\chapter}{0pt}{-12pt}{24pt}
\setlist{topsep=4pt,itemsep=3pt,leftmargin=1.5em}
\captionsetup{font=small,labelfont={bf,color=azulNoche},justification=centering}
\setlength{\parskip}{4pt}
\setlength{\parindent}{1.2em}
\emergencystretch=3em

% --- Enlaces y encabezados ---
\usepackage{hyperref}
\hypersetup{
  colorlinks=true,
  linkcolor=azulNoche,
  urlcolor=azulMedio,
  citecolor=azulMedio,
  pdftitle={Generador Java UML Modelos},
  pdfsubject={Proyecto final de Desarrollo Basado en Modelos},
  pdfauthor={Estudiante}
}
\usepackage{fancyhdr}
\pagestyle{fancy}
\fancyhf{}
\fancyhead[L]{\small\color{grisTexto}Generador Java UML Modelos}
\fancyhead[R]{\small\color{grisTexto}Desarrollo Basado en Modelos}
\fancyfoot[L]{\small\color{grisTexto}Documento académico}
\fancyfoot[R]{\small\color{grisTexto}\thepage}
\renewcommand{\headrulewidth}{0.4pt}
\renewcommand{\footrulewidth}{0pt}
\setlength{\headheight}{16pt}

% --- Código fuente ---
\lstdefinestyle{sql}{
  language=SQL,
  basicstyle=\ttfamily\small,
  keywordstyle=\color{azulMedio}\bfseries,
  commentstyle=\color{grisTexto}\itshape,
  stringstyle=\color{turquesa},
  breaklines=true,
  frame=single,
  rulecolor=\color{azulMedio},
  backgroundcolor=\color{grisFondo},
  showstringspaces=false,
  numbers=left,
  numberstyle=\tiny\color{grisTexto},
  xleftmargin=1em,
  framexleftmargin=0.7em
}

% Insertar capturas reales cuando estén disponibles. Los recuadros alternativos
% mantienen el documento compilable durante la fase de diseño.
\newcommand{\evidencia}[3]{%
  \begin{figure}[H]
    \centering
    \IfFileExists{#1}{\includegraphics[width=0.94\textwidth,height=0.62\textheight,keepaspectratio]{#1}}{%
      \fcolorbox{azulMedio}{grisFondo}{%
        \parbox[c][5.2cm][c]{0.88\textwidth}{\centering
          {\Large\bfseries\color{azulNoche}Evidencia pendiente}\par\vspace{0.7em}
          \textbf{#2}\par\vspace{0.5em}
          {\color{grisTexto}#3}\par\vspace{0.6em}
          {\small\ttfamily\color{azulMedio}\detokenize{#1}}
        }%
      }%
    }
    \caption{#2}
    \label{fig:#1}
  \end{figure}
}

\begin{document}

% ==================== PORTADA ====================
\begin{titlepage}
\thispagestyle{empty}
\begin{tikzpicture}[remember picture,overlay]
  \fill[azulNoche] (current page.north west) rectangle ([yshift=-6.0cm]current page.north east);
  \fill[turquesa] ([yshift=-6.0cm]current page.north west) rectangle ([yshift=-6.18cm]current page.north east);
  \fill[dorado] ([xshift=2.7cm,yshift=-0.75cm]current page.north west) rectangle ([xshift=3.35cm,yshift=-0.84cm]current page.north west);
\end{tikzpicture}
\vspace*{-1.75cm}
\begin{flushleft}
{\large\color{white}\bfseries PROYECTO FINAL}\par
\vspace{0.35cm}
{\LARGE\color{white}\bfseries Generador Java\par UML Modelos}\par
\vspace{0.25cm}
{\large\color{white}Transformación de modelos UML a modelo relacional, Java y SQL}\par
\end{flushleft}
\vspace{5.8cm}
\begin{center}
{\Large\bfseries\color{azulNoche}INFORME TÉCNICO}\par
\vspace{0.6cm}
{\large Asignatura: Desarrollo Basado en Modelos}\par
\end{center}
\vfill
\begin{tabularx}{\textwidth}{>{\bfseries\color{azulNoche}}p{4.0cm}X}
Institución: & \rule{\linewidth}{0.4pt}\\[0.6cm]
Programa académico: & \rule{\linewidth}{0.4pt}\\[0.6cm]
Estudiante: & \rule{\linewidth}{0.4pt}\\[0.6cm]
Docente: & \rule{\linewidth}{0.4pt}\\[0.6cm]
Lugar y fecha: & \rule{\linewidth}{0.4pt}\\
\end{tabularx}
\vspace{1.0cm}
\begin{center}
{\small\color{grisTexto}Versión documental 1.0 \textbar\ Estado: prototipo con validación Java y SQLite}\par
{\small\color{grisTexto}26 de noviembre de 2026}
\end{center}
\end{titlepage}

\pagenumbering{roman}
\chapter*{Resumen}
\addcontentsline{toc}{chapter}{Resumen}
Este proyecto implementa un prototipo de desarrollo basado en modelos que lee un modelo UML serializado en XMI y genera un modelo relacional intermedio en JSON, clases Java 17 y un esquema DDL para SQLite. La ejecución comprobada se realiza con un generador Java autónomo; Eclipse, ATL, Acceleo y Papyrus se prepararon como herramientas de evolución, pero la cadena ATL--Acceleo no se reporta como ejecutada.

El ejemplo incluido contiene Autor, Libro y Editorial, atributos básicos, una relación uno a muchos y una relación muchos a muchos. La ejecución produjo tres entidades y cuatro tablas; el código Java resultante también se compiló con JDK 17. Eclipse abre el proyecto con cero errores y una advertencia en la vista de problemas. El DDL se ejecutó en SQLite 3.54.1: se comprobaron cuatro tablas, relaciones válidas y el rechazo de una clave foránea inválida.

\noindent\textbf{Palabras clave:} desarrollo basado en modelos, UML, EMF, ATL, Acceleo, Java, SQLite.

\tableofcontents
\clearpage
\listoffigures
\clearpage
\listoftables
\clearpage
\pagenumbering{arabic}

% ==================== CAPÍTULO 1 ====================
\chapter{Introducción}
\section{Contexto}
El desarrollo basado en modelos (MDE, por sus siglas en inglés) utiliza modelos como artefactos centrales del proceso de desarrollo. En vez de mantener manualmente la misma estructura en diagramas, clases y bases de datos, una cadena de transformación permite derivar esos artefactos desde una fuente conceptual común.

\section{Problema}
Cuando el diseño conceptual, el código y el esquema de datos se mantienen por separado, un cambio puede propagarse de forma incompleta. Esto produce discrepancias entre nombres, tipos, claves y relaciones. El proyecto busca reducir esa duplicación mediante una transformación reproducible desde UML hacia un modelo relacional y dos salidas textuales.

\section{Objetivo general}
Diseñar e implementar un prototipo MDE que convierta un modelo UML de clases válido en un modelo relacional, clases Java y un script DDL ejecutable en SQLite.

\section{Objetivos específicos}
\begin{enumerate}
  \item Delimitar un subconjunto de UML y acordar las convenciones de tipos, identificadores y nombres.
  \item Definir un metamodelo relacional con EMF/Ecore.
  \item Implementar en ATL la transformación de UML al modelo relacional.
  \item Crear plantillas Acceleo para generar Java y SQL desde el modelo relacional.
  \item Validar las salidas con compilación Java y ejecución del DDL en SQLite.
\end{enumerate}

\section{Enfoque de trabajo}
El trabajo se organizará en incrementos verificables: especificación del subconjunto, metamodelado, transformación, generación, validación y documentación. Las decisiones marcadas como propuestas se confirmarán antes de fijarlas como reglas de implementación.

% ==================== CAPÍTULO 2 ====================
\chapter{Alcance y requisitos}
\section{Alcance funcional}
La versión inicial contempla:
\begin{itemize}
  \item Clases UML persistibles y atributos de tipos básicos.
  \item Asociaciones binarias uno a uno, uno a muchos y muchos a muchos.
  \item Tablas, columnas, claves primarias, claves foráneas y tablas de unión.
  \item Una clase Java por entidad y un script DDL para SQLite.
  \item Diagnósticos para modelos que incumplan las reglas acordadas.
\end{itemize}

\section{Fuera del alcance}
La primera versión no incluye herencia, interfaces, operaciones UML, asociaciones n-arias, composición con semántica especial, aplicación CRUD, interfaz gráfica, sincronización bidireccional ni migraciones de esquemas.

\section{Requisitos funcionales}
\begin{longtable}{p{0.13\textwidth}p{0.79\textwidth}}
\caption{Requisitos funcionales iniciales}\label{tab:rf}\\
\toprule
\textbf{ID} & \textbf{Descripción}\\
\midrule
\endfirsthead
\toprule
\textbf{ID} & \textbf{Descripción}\\
\midrule
\endhead
RF-01 & Leer un modelo UML de clases de Papyrus que cumpla el subconjunto soportado.\\
RF-02 & Validar nombres, tipos, extremos de asociaciones y multiplicidades.\\
RF-03 & Producir un modelo relacional con tablas, columnas y claves.\\
RF-04 & Crear una tabla de unión por cada asociación muchos a muchos.\\
RF-05 & Generar código Java a partir del modelo relacional.\\
RF-06 & Generar DDL para crear el esquema en SQLite.\\
RF-07 & Informar entradas inválidas y evitar presentar salidas incompletas como válidas.\\
\bottomrule
\end{longtable}

\section{Requisitos no funcionales}
Las transformaciones deberán ser reproducibles y deterministas. La organización separará los modelos de entrada, las transformaciones y los artefactos generados. El código generado tendrá como objetivo Java 17; las dependencias y sus versiones se registrarán en la documentación del entorno.

% ==================== CAPÍTULO 3 ====================
\chapter{Arquitectura propuesta}
\section{Cadena de transformación}
La arquitectura introduce un modelo relacional intermedio para hacer explícitas las decisiones de persistencia antes de generar código y SQL. La Figura~\ref{fig:arquitectura} resume el flujo propuesto.

\begin{figure}[H]
\centering
\begin{tikzpicture}[scale=0.85,transform shape,
  node distance=0.75cm,
  bloque/.style={draw=azulNoche,fill=grisFondo,rounded corners=3pt,minimum width=3.0cm,minimum height=1.0cm,align=center,font=\small},
  paso/.style={draw=turquesa,fill=turquesa!8,rounded corners=3pt,minimum width=2.5cm,minimum height=0.85cm,align=center,font=\small},
  flecha/.style={-{Stealth[length=2.5mm]},thick,draw=azulMedio}
]
\node[bloque] (uml) {Modelo UML\\Papyrus};
\node[paso,right=of uml] (atl) {Validación\\y ATL};
\node[bloque,right=of atl] (rel) {Modelo relacional\\EMF/Ecore};
\node[paso,right=of rel] (acc) {Plantillas\\Acceleo};
\node[bloque,above right=0.25cm and 0.55cm of acc] (java) {Clases Java\\Java 17};
\node[bloque,below right=0.25cm and 0.55cm of acc] (sql) {Esquema DDL\\SQLite};
\draw[flecha] (uml) -- (atl);
\draw[flecha] (atl) -- (rel);
\draw[flecha] (rel) -- (acc);
\draw[flecha] (acc.east) -- (java.west);
\draw[flecha] (acc.east) -- (sql.west);
\end{tikzpicture}
\caption{Arquitectura lógica propuesta para el flujo MDE.}
\label{fig:arquitectura}
\end{figure}

\section{Responsabilidades de los componentes}
\begin{longtable}{p{0.22\textwidth}p{0.68\textwidth}}
\caption{Responsabilidades por componente}\label{tab:componentes}\\
\toprule
\textbf{Componente} & \textbf{Responsabilidad}\\
\midrule
\endfirsthead
\toprule
\textbf{Componente} & \textbf{Responsabilidad}\\
\midrule
\endhead
Papyrus/UML & Edición del modelo conceptual de clases y asociaciones.\\
EMF/Ecore & Definición de metamodelos y serialización de modelos.\\
ATL & Transformación de elementos UML admitidos al modelo relacional.\\
Acceleo & Generación de archivos Java y SQL a partir del modelo relacional.\\
Java & Compilación y posible extensión de las clases generadas; objetivo inicial Java 17.\\
SQLite & Ejecución y comprobación del esquema generado.\\
\bottomrule
\end{longtable}

\section{Principios de diseño}
El modelo UML será la fuente conceptual. El modelo relacional será la fuente de generación. ATL contendrá las reglas de persistencia y Acceleo se ocupará de presentar el modelo relacional como texto. Los archivos generados serán reproducibles y no deberán editarse manualmente. Para modelos fuera del alcance se emitirá un diagnóstico explícito.

% ==================== CAPÍTULO 4 ====================
\chapter{Reglas de generación implementadas}
\section{Clases, atributos e identificadores}
Cada clase UML persistible se transforma en una entidad relacional, una tabla y una clase Java. Cada atributo admitido se convierte en columna y propiedad Java.

\begin{quote}
Si una clase no declara un identificador, el generador crea una clave técnica entera \texttt{id}, utilizada como clave primaria en el esquema SQL.
\end{quote}

\section{Correspondencia de tipos}
La tabla~\ref{tab:tipos} documenta la correspondencia implementada por el prototipo.

\begin{table}[H]
\centering
\caption{Correspondencia inicial de tipos}
\label{tab:tipos}
\begin{tabular}{lll}
\toprule
\textbf{Tipo UML} & \textbf{Tipo Java sugerido} & \textbf{Tipo SQLite sugerido}\\
\midrule
String & String & TEXT\\
Integer & Integer & INTEGER\\
Boolean & Boolean & INTEGER (0/1)\\
Real & Double & REAL\\
\bottomrule
\end{tabular}
\end{table}

Los atributos obligatorios se generan como \texttt{NOT NULL}; los booleanos se representan como enteros con una restricción a 0 o 1.

\section{Asociaciones}
\begin{table}[H]
\centering
\caption{Representación relacional propuesta para asociaciones}
\label{tab:asociaciones}
\begin{tabularx}{\textwidth}{l l X}
\toprule
\textbf{Relación UML} & \textbf{Representación} & \textbf{Restricción propuesta}\\
\midrule
Uno a uno & Clave foránea en una tabla & La clave foránea tendrá restricción de unicidad.\\
Uno a muchos & Clave foránea en el lado muchos & Referencia la clave primaria del lado uno.\\
Muchos a muchos & Tabla de unión & Dos claves foráneas con clave primaria compuesta.\\
\bottomrule
\end{tabularx}
\end{table}

Las tablas de unión usan una convención determinista basada en los nombres de las clases; su clave primaria es compuesta por las dos claves foráneas.

\section{Ejemplo de salida SQL esperada}
El siguiente fragmento es parte del DDL generado para las entidades \texttt{Autor} y \texttt{Libro} en el ejemplo incluido.

\begin{lstlisting}[style=sql,caption={DDL ilustrativo para una relación muchos a muchos}]
CREATE TABLE Autor (
  id INTEGER PRIMARY KEY,
  nombre TEXT NOT NULL
);

CREATE TABLE Libro (
  id INTEGER PRIMARY KEY,
  titulo TEXT NOT NULL
);

CREATE TABLE Autor_Libro (
  autor_id INTEGER NOT NULL,
  libro_id INTEGER NOT NULL,
  PRIMARY KEY (autor_id, libro_id),
  FOREIGN KEY (autor_id) REFERENCES Autor(id),
  FOREIGN KEY (libro_id) REFERENCES Libro(id)
);
\end{lstlisting}

% ==================== CAPÍTULO 5 ====================
\chapter{Validación y manejo de errores}
El prototipo comprueba las condiciones necesarias para generar el modelo de ejemplo. Su cobertura no equivale a una validación completa del estándar UML. Entre las condiciones que revisa se encuentran:
\begin{enumerate}
  \item Los nombres de clases y atributos sean válidos y únicos en su ámbito.
  \item Los tipos de atributos pertenezcan al catálogo admitido.
  \item Las asociaciones sean binarias y sus extremos estén correctamente definidos.
  \item Las multiplicidades puedan representarse con las reglas implementadas.
  \item Los nombres resultantes de tablas, columnas y restricciones no colisionen.
\end{enumerate}

Los tipos admitidos son String, Integer, Boolean y Real/Double. Las asociaciones se limitan al subconjunto binario procesado por el lector. Los mensajes ante entradas fuera de alcance se reportan durante la ejecución.

% ==================== CAPÍTULO 6 ====================
\chapter{Entorno y organización del repositorio}
\section{Herramientas previstas}
El entorno de trabajo preparado utiliza Eclipse Modeling Tools 2025-06, Git y Java 17 para las salidas generadas. La distribución de Eclipse incorpora una JVM para ejecutar el IDE. El entorno incluye Papyrus Desktop, ATL y Acceleo desde sitios oficiales. Eclipse abre el espacio de trabajo y el proyecto sin errores, con una advertencia visible. El flujo ATL/Acceleo no forma parte de la ejecución verificada. Las versiones y fuentes se documentan en \texttt{docs/TOOLCHAIN.md}.

\begin{itemize}
  \item Papyrus Desktop 7.1.0 se publica para Eclipse 2025-06.
  \item ATL 4.12.0 está disponible desde el repositorio oficial de versiones.
  \item Acceleo 3.7.17 quedó registrado en la instalación; no se ha validado una transformación ATL/Acceleo ejecutable.
\end{itemize}

SQLite es embebido y no requiere un servidor. SQLite 3.54.1 se instaló localmente en E: y el script de ejecución valida el DDL en una base temporal, comprueba cuatro tablas y prueba restricciones referenciales.

\section{Estructura del repositorio y extensiones previstas}
\begin{description}[leftmargin=3.5cm,style=nextline]
  \item[\texttt{models/uml/}] Modelos Papyrus de entrada.
  \item[\texttt{models/relational/}] Metamodelo Ecore y modelos relacionales.
  \item[\texttt{transformations/atl/}] Reglas de UML a relacional.
  \item[\texttt{generator/acceleo/}] Plantillas para Java y SQL.
  \item[\texttt{examples/biblioteca/}] Ejemplo para recorrer el flujo.
  \item[\texttt{generated/java/}] Clases Java generadas.
  \item[\texttt{generated/sql/}] Scripts SQL generados.
  \item[\texttt{docs/}] Documentación y guías del entorno.
\end{description}

\section{Procedimiento de uso previsto}
\begin{enumerate}
  \item Crear o editar el modelo UML en Papyrus.
  \item Validar que cumpla el subconjunto soportado.
  \item Ejecutar ATL y revisar el modelo relacional resultante.
  \item Ejecutar Acceleo para generar Java y SQL.
  \item Ejecutar el DDL en SQLite y comprobar tablas y restricciones.
  \item Regenerar después de cada cambio UML, sin editar manualmente las salidas.
\end{enumerate}

% ==================== CAPÍTULO 7 ====================
\chapter{Estado del proyecto y evidencias}
\section{Estado verificable}
El prototipo se ejecutó desde E:. Se compiló el generador y sus tres clases producidas con JDK 17. SQLite 3.54.1 creó cuatro tablas desde el DDL, aceptó relaciones válidas y rechazó una clave foránea inválida. Eclipse abrió el proyecto con cero errores y una advertencia. ATL y Acceleo siguen sin integrarse en el ejecutable verificado.

\section{Evidencia del entorno de trabajo}
Eclipse abre el proyecto, pero no se pudo guardar una captura durante esta sesión porque la superposición de la Herramienta Recortes no respondió a las acciones de captura.

\evidencia{figuras/eclipse-proyecto.png}
{Proyecto abierto en Eclipse}
{Pendiente: insertar una captura real de Eclipse; no se presenta una imagen simulada.}

\section{Evidencia de ejecución}
La ejecución de scripts/run.ps1 se verificó, pero la captura no pudo guardarse desde la herramienta de Recortes en esta sesión.

\evidencia{figuras/ejecucion-generador.png}
{Ejecución del generador y producción de salidas}
{Pendiente: insertar una captura auténtica de consola y archivos generados; no se presenta una imagen simulada.}

\section{Resultados de validación}
\begin{table}[H]
\centering
\caption{Resultados de comprobación del prototipo}
\label{tab:pruebas}
\begin{tabularx}{\textwidth}{l X X}
\toprule
\textbf{ID} & \textbf{Comprobación} & \textbf{Resultado}\\
\midrule
PV-01 & Lectura de las tres clases del modelo & Correcto: tres entidades.\\
PV-02 & Generación de Java y modelo intermedio & Correcto: tres clases y JSON.\\
PV-03 & Compilación del Java generado & Correcto con JDK 17.\\
PV-04 & Asociación muchos a muchos & Correcto: tabla unión y claves foráneas.\\
PV-05 & Asociación uno a muchos & Correcto: FK de Libro a Editorial.\\
PV-06 & Ejecución del DDL en SQLite & Correcto: cuatro tablas; relaciones y FK inválida comprobadas.\\
PV-07 & Apertura del proyecto en Eclipse & Correcto: cero errores y una advertencia.\\
\bottomrule
\end{tabularx}
\end{table}

Se verificaron lectura, generación, compilación Java, DDL en SQLite y apertura del proyecto en Eclipse. La captura real de las pantallas queda pendiente; también siguen pendientes la integración ATL/Acceleo y el uso de Papyrus para editar el modelo.

% ==================== CAPÍTULO 8 ====================
\chapter{Plan de trabajo}
\section{Hitos}
\begin{enumerate}
  \item Abrir y editar el modelo de ejemplo en Papyrus.
  \item Definir el metamodelo relacional Ecore y ejecutar una transformación ATL.
  \item Implementar y ejecutar las plantillas Acceleo.
  \item Integrar una transformación ATL y plantillas Acceleo en la ejecución reproducible.
  \item Guardar capturas reales del proyecto en Eclipse, la ejecución Java/SQLite y los archivos generados.
  \item Ampliar las reglas y validar entradas inválidas.
  \item Documentar los resultados medidos y limitaciones.
\end{enumerate}

\section{Riesgos y mitigación}
UML admite semánticas más amplias que el subconjunto inicial. Para limitar errores, el alcance deberá validarse de forma temprana. Identificadores, nombres reservados y colisiones requieren reglas explícitas. La combinación de Eclipse y complementos debe mantenerse documentada. Por último, el generador deberá validar referencias antes de escribir salidas parciales.

\chapter{Conclusiones}
Se implementó un prototipo reproducible que transforma una entrada UML/XMI pequeña en un modelo relacional JSON, clases Java y DDL SQL. La generación y compilación Java se verificaron con JDK 17. SQLite 3.54.1 ejecutó el DDL y confirmó las cuatro tablas, relaciones válidas y rechazo de una clave foránea inválida. Eclipse abre el proyecto con cero errores y una advertencia. La integración ATL/Acceleo, el uso de Papyrus y la incorporación de capturas reales siguen como trabajo pendiente.
\chapter*{Glosario}
\addcontentsline{toc}{chapter}{Glosario}
\begin{description}[leftmargin=3cm,style=nextline]
  \item[Acceleo] Generador de texto basado en modelos EMF y MTL.
  \item[ATL] Lenguaje y entorno para transformar modelos.
  \item[DDL] Sentencias SQL que definen objetos de una base de datos.
  \item[EMF/Ecore] Marco de Eclipse para definir metamodelos y modelos.
  \item[MDE] Desarrollo de software basado en modelos.
  \item[Modelo fuente] Modelo UML inicial del flujo.
  \item[Modelo destino] Modelo relacional resultante.
\end{description}

\chapter*{Referencias}
\addcontentsline{toc}{chapter}{Referencias}
\begin{enumerate}
  \item Eclipse Foundation, \emph{Papyrus Downloads}. \url{https://eclipse.dev/papyrus/download.html}
  \item Eclipse Foundation, \emph{ATL release repository}. \url{https://download.eclipse.org/mmt/atl/builds/release/}
  \item Eclipse Foundation, \emph{Acceleo 3.7.16}. \url{https://projects.eclipse.org/projects/modeling.acceleo/releases/3.7.16}
  \item SQLite Project, \emph{SQLite Documentation}. \url{https://www.sqlite.org/docs.html}
  \item Eclipse Foundation, \emph{Eclipse Modeling Tools 2025-06}. \url{https://www.eclipse.org/downloads/packages/release/2025-06/r/eclipse-modeling-tools}
\end{enumerate}

\end{document}
